\documentclass[aspectratio=169, 10pt]{beamer}

% ============================================================
%  TEMA Y COLORES
% ============================================================
\usetheme{CambridgeUS}
\usecolortheme{seahorse}

\definecolor{fenblue}{RGB}{0,56,101}
\definecolor{fenred}{RGB}{180,30,45}
\definecolor{fengray}{RGB}{90,90,90}
\definecolor{fenlight}{RGB}{235,241,248}

\setbeamercolor{frametitle}{bg=fenblue, fg=white}
\setbeamercolor{title}{fg=fenblue}
\setbeamercolor{structure}{fg=fenblue}
\setbeamercolor{block title}{bg=fenblue, fg=white}
\setbeamercolor{block body}{bg=fenlight}
\setbeamercolor{alerted text}{fg=fenred}
\setbeamercolor{itemize item}{fg=fenblue}
\setbeamercolor{itemize subitem}{fg=fengray}
\setbeamercolor{footline}{bg=fenblue, fg=white}

\setbeamertemplate{navigation symbols}{}
\setbeamertemplate{footline}{%
  \leavevmode%
  \hbox{%
    \begin{beamercolorbox}[wd=.6\paperwidth,ht=2.25ex,dp=1ex,leftskip=.3cm]{footline}%
      \usebeamerfont{author in head/foot}\insertshortauthor \quad | \quad \insertshorttitle
    \end{beamercolorbox}%
    \begin{beamercolorbox}[wd=.4\paperwidth,ht=2.25ex,dp=1ex,rightskip=.3cm plus1fil]{footline}%
      \usebeamerfont{date in head/foot}\hfill\insertshortdate\quad|\quad\insertframenumber/\inserttotalframenumber
    \end{beamercolorbox}}%
  \vskip0pt%
}

% ============================================================
%  PAQUETES
% ============================================================
\usepackage[spanish]{babel}
\usepackage[utf8]{inputenc}
\usepackage[T1]{fontenc}
\usepackage{lmodern}
\usepackage{tikz}
\usepackage{booktabs}
\usepackage{array}
\usepackage{xcolor}
\usepackage{fontawesome5}

\usetikzlibrary{arrows.meta, positioning, shapes.geometric, fit, backgrounds}

% ============================================================
%  METADATOS
% ============================================================
\title[Attention Is All You Need]{\textbf{Attention Is All You Need}}
\subtitle{Vaswani et al. (2017) --- El paper que cambió el NLP}
\author[Sebastián Egaña]{Sebastián Egaña Santibáñez}
\institute[FEN Unegocios]{Módulo 3 · IA Generativa y LLMs \\ FEN Unegocios, Universidad de Chile}
\date[2026]{Abril 2026}

% ============================================================
%  MACROS UTILITARIOS
% ============================================================
\newcommand{\highlight}[1]{\textcolor{fenred}{\textbf{#1}}}
\newcommand{\concept}[1]{\textcolor{fenblue}{\textbf{#1}}}

% ============================================================
\begin{document}
% ============================================================

% ----------------------------------------------------------
%  PORTADA
% ----------------------------------------------------------
\begin{frame}[plain]
  \maketitle
\end{frame}

% ----------------------------------------------------------
%  AGENDA
% ----------------------------------------------------------
\begin{frame}{Agenda}
  \tableofcontents[hideallsubsections]
\end{frame}

% ===========================================================
\section{El problema que existía}
% ===========================================================

\begin{frame}{Antes de 2017: modelos que leían de izquierda a derecha}

  \begin{columns}[T]
    \begin{column}{0.55\textwidth}
      \textbf{Modelo predominante: RNN / LSTM}
      \vspace{0.3em}
      \begin{itemize}
        \item Procesan texto \concept{token a token}, en secuencia
        \item Para entender la palabra 10, deben ``recordar'' las 9 anteriores
        \item La información tiene que \highlight{sobrevivir} muchos pasos intermedios
      \end{itemize}
      \vspace{0.5em}
      \begin{alertblock}{El límite práctico}
        En frases largas, el modelo \textit{olvidaba} detalles del principio antes de llegar al final.
      \end{alertblock}
    \end{column}
    \begin{column}{0.42\textwidth}
      \begin{tikzpicture}[
        node distance=0.6cm,
        token/.style={draw=fenblue, rounded corners, fill=fenlight,
                      minimum width=1.2cm, minimum height=0.6cm, font=\small},
        arr/.style={-{Latex}, thick, fenblue}
      ]
        \node[token] (t1) {el};
        \node[token, below=of t1] (t2) {modelo};
        \node[token, below=of t2] (t3) {ley\'o};
        \node[token, below=of t3] (t4) {el};
        \node[token, below=of t4] (t5) {informe};
        \draw[arr] (t1)--(t2);
        \draw[arr] (t2)--(t3);
        \draw[arr] (t3)--(t4);
        \draw[arr] (t4)--(t5);
        \node[right=0.3cm of t1, font=\tiny, fengray] {estado};
        \node[right=0.3cm of t2, font=\tiny, fengray] {estado};
        \node[right=0.3cm of t3, font=\tiny, fengray] {estado};
        \node[right=0.3cm of t4, font=\tiny, fengray] {estado};
        \node[right=0.3cm of t5, font=\tiny, fengray] {estado};
      \end{tikzpicture}
      \vspace{0.2em}
      {\small \textit{Cada flecha pasa información al siguiente paso}}
    \end{column}
  \end{columns}
\end{frame}

% ----------------------------------------------------------

\begin{frame}{El cuello de botella: dependencias largas}

  \begin{block}{Ejemplo del curso}
    \textit{``El analista revisó el informe porque necesitaba verificar si los \textbf{beneficios} estaban correctamente \textbf{calculados}.''}
  \end{block}

  \vspace{0.5em}
  \begin{columns}[T]
    \begin{column}{0.48\textwidth}
      \textbf{Problema con RNN:}
      \begin{itemize}
        \item Para conectar \textit{beneficios} con \textit{calculados}, la información debe pasar por todos los tokens intermedios
        \item Cuantos más tokens intermedios $\rightarrow$ más probabilidad de ``diluirse''
        \item No puede paralelizarse $\rightarrow$ entrenamiento lento
      \end{itemize}
    \end{column}
    \begin{column}{0.48\textwidth}
      \textbf{¿Qué querríamos?}
      \begin{itemize}
        \item Que cualquier palabra pudiera ``hablarle directamente'' a cualquier otra
        \item Sin importar la distancia
        \item \concept{Y en paralelo}, para escalar
      \end{itemize}
    \end{column}
  \end{columns}

  \vspace{0.6em}
  \begin{center}
    \large \textbf{Esa es exactamente la propuesta del paper.}
  \end{center}

\end{frame}

% ===========================================================
\section{La propuesta: solo atención}
% ===========================================================

\begin{frame}{``Attention is all you need'' --- la apuesta central}

  \begin{columns}[T]
    \begin{column}{0.55\textwidth}
      \begin{block}{La tesis del paper (Vaswani et al., 2017)}
        Es posible construir un modelo de lenguaje de alta calidad \highlight{sin recurrencia ni convoluciones}, usando únicamente mecanismos de atención.
      \end{block}

      \vspace{0.5em}
      \textbf{Lo que propone el Transformer:}
      \begin{enumerate}
        \item Cada token mira a \textit{todos} los demás tokens a la vez
        \item Les asigna un peso según cuán relevantes son
        \item Construye su representación a partir de ese peso
        \item \concept{Todo en paralelo} --- no en secuencia
      \end{enumerate}
    \end{column}
    \begin{column}{0.42\textwidth}
      \begin{tikzpicture}[
        token/.style={draw=fenblue, rounded corners, fill=fenlight,
                      minimum width=1cm, minimum height=0.5cm, font=\scriptsize},
        arr/.style={-{Latex}, fenred, thick}
      ]
        \node[token] (A) at (0,2.5) {el};
        \node[token] (B) at (1.8,2.5) {modelo};
        \node[token] (C) at (0,0.8) {ley\'o};
        \node[token] (D) at (1.8,0.8) {informe};

        \draw[arr] (A) to[bend left=15] (B);
        \draw[arr] (A) to[bend right=15] (C);
        \draw[arr] (A) to[bend left=10] (D);
        \draw[arr] (B) to[bend left=15] (C);
        \draw[arr] (B) to[bend right=10] (D);
        \draw[arr] (C) to[bend left=15] (D);

        \node[below=0.3cm of C, font=\tiny, align=center, fengray] {Cada token\\se conecta con todos};
      \end{tikzpicture}
    \end{column}
  \end{columns}

\end{frame}

% ===========================================================
\section{El mecanismo de atención}
% ===========================================================

\begin{frame}{La intuición: ¿qué es ``prestar atención''?}

  \begin{block}{Definición operativa}
    La atención es el mecanismo que permite al modelo \concept{ponderar qué partes del contexto son más relevantes} en cada momento.
  \end{block}

  \vspace{0.5em}
  \begin{columns}[T]
    \begin{column}{0.55\textwidth}
      Cuando el modelo genera o representa una palabra:
      \begin{itemize}
        \item Mira \textit{todos} los tokens del contexto
        \item Estima cuáles son más relevantes
        \item Les da \highlight{más peso} a algunos
        \item Y \highlight{menos peso} a otros
      \end{itemize}
      \vspace{0.4em}
      \textbf{No recuerda todo igual: prioriza dinámicamente.}
    \end{column}
    \begin{column}{0.42\textwidth}
      \begin{exampleblock}{Ejemplo: ``banco''}
        \textit{``Fui al \textbf{banco} a depositar el cheque.''}\\[0.3em]
        Tokens de alto peso: \texttt{depositar}, \texttt{cheque}\\[0.2em]
        Tokens de bajo peso: \texttt{Fui}, \texttt{al}, \texttt{a}\\[0.3em]
        $\Rightarrow$ El modelo infiere ``banco'' = institución financiera
      \end{exampleblock}
    \end{column}
  \end{columns}

\end{frame}

% ----------------------------------------------------------

\begin{frame}{Queries, Keys y Values --- sin fórmulas}

  \begin{columns}[T]
    \begin{column}{0.53\textwidth}
      Cada token genera \textbf{tres vectores}:
      \vspace{0.4em}
      \begin{description}
        \item[\concept{Query (Q)}] \textit{¿Qué estoy buscando?}\\
          La pregunta que hace el token actual
        \item[\concept{Key (K)}] \textit{¿Qué tipo de información tengo?}\\
          Lo que cada token ``ofrece''
        \item[\concept{Value (V)}] \textit{¿Cuál es mi contenido?}\\
          La información que entregará si es relevante
      \end{description}
      \vspace{0.4em}
      El modelo compara Q con cada K $\rightarrow$ obtiene un puntaje $\rightarrow$ pondera los V.
    \end{column}
    \begin{column}{0.44\textwidth}
      \begin{exampleblock}{Analogía: pregunta sobre pensiones}
        \textbf{Pregunta (Query):}\\
        ``¿Qué documentos necesito para pensionarme?''\\[0.3em]
        \textbf{Fragmentos relevantes (Keys con alto match):}\\
        {\small certificados, edad legal, formulario, solicitud}\\[0.3em]
        \textbf{Contenido que se ponderará más (Values):}\\
        {\small requisitos, plazos, documentación}
      \end{exampleblock}
    \end{column}
  \end{columns}

\end{frame}

% ----------------------------------------------------------

\begin{frame}{Scaled Dot-Product Attention --- la mecánica}

  \begin{columns}[T]
    \begin{column}{0.50\textwidth}
      \textbf{En palabras simples:}
      \begin{enumerate}
        \item Para el token actual, toma su \concept{Query}
        \item Compárala contra las \concept{Keys} de todos los demás tokens (producto punto)
        \item Normaliza los puntajes (softmax $\Rightarrow$ suman 1)
        \item Multiplica cada puntaje por el \concept{Value} correspondiente
        \item Suma ponderada $=$ nueva representación del token
      \end{enumerate}

      \vspace{0.3em}
      \begin{alertblock}{Por qué ``scaled''}
        El producto punto crece con la dimensión del vector. Se divide por $\sqrt{d_k}$ para estabilizar el gradiente durante el entrenamiento.
      \end{alertblock}
    \end{column}
    \begin{column}{0.47\textwidth}
      \begin{center}
      \begin{tikzpicture}[
        box/.style={draw=fenblue, fill=fenlight, rounded corners,
                    minimum width=2.2cm, minimum height=0.55cm,
                    font=\scriptsize, align=center},
        arr/.style={-{Latex}, thick, fenblue}
      ]
        \node[box] (qkv) {Matrices Q, K, V};
        \node[box, below=0.5cm of qkv] (dot) {Producto punto QK$^T$};
        \node[box, below=0.5cm of dot] (scale) {Escalar ($\div\sqrt{d_k}$)};
        \node[box, below=0.5cm of scale] (soft) {Softmax};
        \node[box, below=0.5cm of soft] (val) {Ponderar por V};
        \node[box, below=0.5cm of val, fill=fenblue!20] (out) {\textbf{Salida de atenci\'on}};

        \draw[arr] (qkv)--(dot);
        \draw[arr] (dot)--(scale);
        \draw[arr] (scale)--(soft);
        \draw[arr] (soft)--(val);
        \draw[arr] (val)--(out);
      \end{tikzpicture}
      \end{center}
    \end{column}
  \end{columns}

\end{frame}

% ----------------------------------------------------------

\begin{frame}{Multi-Head Attention --- mirar con varios lentes}

  \begin{columns}[T]
    \begin{column}{0.54\textwidth}
      \textbf{¿Por qué múltiples cabezas?}
      \begin{itemize}
        \item Una sola atención captura \textit{un tipo} de relación a la vez
        \item Con múltiples cabezas en paralelo, el modelo puede capturar \concept{distintos tipos de dependencias simultáneamente}:
          \begin{itemize}
            \item {\small sintácticas (sujeto-verbo)}
            \item {\small semánticas (causa-efecto)}
            \item {\small correferenciales (pronombre-referente)}
          \end{itemize}
      \end{itemize}
      \vspace{0.3em}
      Los resultados de cada cabeza se concatenan y se proyectan en una representación final.
    \end{column}
    \begin{column}{0.43\textwidth}
      \begin{center}
      \begin{tikzpicture}[
        head/.style={draw=fenblue, fill=fenlight, rounded corners,
                     minimum width=1.5cm, minimum height=0.5cm,
                     font=\scriptsize},
        arr/.style={-{Latex}, thick, fenblue}
      ]
        \node[head] (in) {Entrada (Q, K, V)};
        \node[head, below left=0.6cm and 0.5cm of in]  (h1) {Head 1};
        \node[head, below=0.8cm of in]                  (h2) {Head 2};
        \node[head, below right=0.6cm and 0.5cm of in] (h3) {Head h};
        \node[head, below=1.6cm of in, fill=fenblue!20]  (cat) {Concat + Proyecci\'on};
        \node[head, below=0.5cm of cat]                 (out) {\textbf{Salida MHA}};

        \draw[arr] (in)--(h1); \draw[arr] (in)--(h2); \draw[arr] (in)--(h3);
        \draw[arr] (h1)--(cat); \draw[arr] (h2)--(cat); \draw[arr] (h3)--(cat);
        \draw[arr] (cat)--(out);

        \node[font=\scriptsize, fengray] at (2.2, -0.8) {$\cdots$};
      \end{tikzpicture}
      \end{center}
    \end{column}
  \end{columns}

\end{frame}

% ===========================================================
\section{La arquitectura Transformer}
% ===========================================================

\begin{frame}{Arquitectura Encoder--Decoder}

  \begin{columns}[T]
    \begin{column}{0.48\textwidth}
      \textbf{Encoder} (comprende la entrada):
      \begin{enumerate}
        \item Input Embeddings + Positional Encoding
        \item Bloque $\times N$:
          \begin{itemize}
            \item \concept{Multi-Head Self-Attention}
            \item Add \& Norm
            \item Feed-Forward Network
            \item Add \& Norm
          \end{itemize}
      \end{enumerate}
    \end{column}
    \begin{column}{0.48\textwidth}
      \textbf{Decoder} (genera la salida):
      \begin{enumerate}
        \item Output Embeddings + Positional Encoding
        \item Bloque $\times N$:
          \begin{itemize}
            \item \concept{Masked} Multi-Head Self-Attention
            \item Add \& Norm
            \item \concept{Cross-Attention} (mira al encoder)
            \item Add \& Norm
            \item Feed-Forward Network
            \item Add \& Norm
          \end{itemize}
        \item Linear + Softmax $\rightarrow$ token de salida
      \end{enumerate}
    \end{column}
  \end{columns}

  \vspace{0.5em}
  \begin{alertblock}{Positional Encoding}
    Como la atención no es secuencial, hay que \textit{inyectar} la posición de cada token. Se suman señales sinusoidales a los embeddings.
  \end{alertblock}

\end{frame}

% ----------------------------------------------------------

\begin{frame}{¿Por qué Masked Self-Attention en el Decoder?}

  \begin{columns}[T]
    \begin{column}{0.55\textwidth}
      \textbf{Problema:} al \textit{generar} texto, el modelo no puede ver palabras que aún no ha producido.

      \vspace{0.4em}
      \textbf{Solución:} durante el entrenamiento se aplica una \concept{máscara} que bloquea el acceso a posiciones futuras.

      \begin{exampleblock}{Intuición}
        Si el modelo ya generó \textit{``El informe fue''}, al predecir la siguiente palabra solo puede mirar esos tres tokens --- no los que vendrán después.
      \end{exampleblock}
    \end{column}
    \begin{column}{0.42\textwidth}
      \begin{center}
      \begin{tikzpicture}[scale=0.8]
        \foreach \i/\w in {1/El, 2/informe, 3/fue, 4/?}{
          \node[draw=fenblue, fill=fenlight, rounded corners,
                minimum width=0.9cm, minimum height=0.5cm,
                font=\scriptsize] (t\i) at (\i*1.1, 0) {\w};
        }
        % allowed
        \draw[fenblue, thick, ->] (t1) to[bend left=30] (t2);
        \draw[fenblue, thick, ->] (t1) to[bend left=40] (t3);
        \draw[fenblue, thick, ->] (t2) to[bend left=30] (t3);
        \draw[fenblue, thick, ->] (t1) to[bend left=50] (t4);
        \draw[fenblue, thick, ->] (t2) to[bend left=40] (t4);
        \draw[fenblue, thick, ->] (t3) to[bend left=30] (t4);
        % masked (show as X)
        \draw[fenred, very thick, ->] (t4) to[bend right=40] node[midway, above, font=\tiny, fenred] {$\times$} (t3);
      \end{tikzpicture}
      \end{center}
      {\small \textit{Solo mira hacia atrás}}
    \end{column}
  \end{columns}

\end{frame}

% ===========================================================
\section{Por qué fue tan influyente}
% ===========================================================

\begin{frame}{El impacto: RNN vs.\ Transformer}

  \begin{center}
  \begin{tabular}{lcc}
    \toprule
    \textbf{Criterio} & \textbf{RNN / LSTM} & \textbf{Transformer} \\
    \midrule
    Procesamiento & Secuencial & \concept{Paralelo} \\
    Dependencias largas & Se diluyen & \concept{Acceso directo} \\
    Escalabilidad & Limitada & \concept{Alta} \\
    Velocidad de entrenamiento & Lenta & \concept{Mucho más rápida} \\
    Captura de contexto & Local & \concept{Global} \\
    \bottomrule
  \end{tabular}
  \end{center}

  \vspace{0.5em}
  \begin{block}{Por qué la paralelización importa tanto}
    Los modelos modernos tienen miles de millones de parámetros. Sin paralelización no sería posible entrenarlos en tiempo razonable. El Transformer es lo que hace viable el escalado a GPT-4, Claude, Gemini, etc.
  \end{block}

\end{frame}

% ----------------------------------------------------------

\begin{frame}{Resultados del paper --- BLEU en traducción}

  \begin{columns}[T]
    \begin{column}{0.5\textwidth}
      \textbf{Tarea: traducción EN$\to$DE e EN$\to$FR (WMT 2014)}

      \vspace{0.4em}
      \begin{tabular}{lc}
        \toprule
        \textbf{Modelo} & \textbf{BLEU EN$\to$DE} \\
        \midrule
        Mejor RNN ensemble & 26.4 \\
        ConvS2S & 26.4 \\
        \concept{Transformer (base)} & \concept{27.3} \\
        \highlight{Transformer (big)} & \highlight{28.4} \\
        \bottomrule
      \end{tabular}

      \vspace{0.4em}
      {\small BLEU = métrica estándar de calidad de traducción (0--100)}
    \end{column}
    \begin{column}{0.46\textwidth}
      \begin{alertblock}{El dato clave}
        Superó al estado del arte con \textbf{menos costo computacional} de entrenamiento que los modelos previos.\\[0.3em]
        No solo mejor calidad: también más eficiente.
      \end{alertblock}
      \vspace{0.3em}
      El paper también probó análisis de constituyentes del inglés, mostrando que la arquitectura generaliza más allá de la traducción.
    \end{column}
  \end{columns}

\end{frame}

% ----------------------------------------------------------

\begin{frame}{De 2017 a hoy: la línea directa}

  \begin{center}
  \begin{tikzpicture}[
    ev/.style={draw=fenblue, rounded corners, fill=fenlight,
               minimum width=2.8cm, minimum height=0.7cm,
               font=\small, align=center},
    arr/.style={-{Latex}, thick, fenblue}
  ]
    \node[ev] (tr)   at (0,0)    {Transformer\\{\tiny Vaswani et al., 2017}};
    \node[ev] (bert) at (3.5,0)  {BERT\\{\tiny Google, 2018}};
    \node[ev] (gpt2) at (7,0)    {GPT-2/3\\{\tiny OpenAI, 2019-20}};
    \node[ev] (ch)   at (3.5,-1.5) {ChatGPT\\{\tiny OpenAI, 2022}};
    \node[ev] (cl)   at (0,-1.5)   {Claude / Gemini\\{\tiny Anthropic/Google, 2023+}};
    \node[ev] (ll)   at (7,-1.5)   {Llama / Mistral\\{\tiny Meta/Mistral, 2023+}};

    \draw[arr] (tr)--(bert);
    \draw[arr] (bert)--(gpt2);
    \draw[arr] (gpt2)--(ch);
    \draw[arr] (tr) to[bend right=20] (cl);
    \draw[arr] (gpt2) to[bend left=20] (ll);
  \end{tikzpicture}
  \end{center}

  \vspace{0.3em}
  \begin{center}
    \large \textbf{Todos son Transformers. La arquitectura de 2017 sigue vigente.}
  \end{center}

\end{frame}

% ===========================================================
\section{Implicancias prácticas}
% ===========================================================

\begin{frame}{Atención y calidad del prompt}

  \begin{columns}[T]
    \begin{column}{0.55\textwidth}
      \textbf{La atención no es una caja negra irrelevante para nosotros.}
      \vspace{0.3em}
      \begin{itemize}
        \item La forma en que escribimos un prompt \concept{afecta qué partes del contexto destacan más}
        \item Si el contexto está mal armado, el modelo puede enfocarse en lo menos importante
        \item Un buen prompt también \textit{guía} la atención del modelo
      \end{itemize}
    \end{column}
    \begin{column}{0.42\textwidth}
      \begin{block}{Por eso ayuda\ldots}
        \begin{itemize}
          \item Separar instrucciones de contexto
          \item Usar formato claro (listas, secciones)
          \item Destacar restricciones críticas
          \item Ordenar bien la información
          \item No enterrar restricciones importantes en texto largo
        \end{itemize}
      \end{block}
    \end{column}
  \end{columns}

\end{frame}

% ----------------------------------------------------------

\begin{frame}{Caso aplicado: instrucción crítica en el agente AFP}

  \begin{block}{Escenario}
    Una AFP implementa un agente que responde preguntas de afiliados. El contexto tiene 5.000 tokens de normativa mezclada.
  \end{block}

  \vspace{0.4em}
  \begin{columns}[T]
    \begin{column}{0.48\textwidth}
      {\color{fenred}\textbf{Prompt mal estructurado:}}
      \vspace{0.2em}

      \begin{quote}\small
        \textit{[\ldots 3.000 tokens de normativa\ldots]}\\
        Importante: si no tienes certeza, di ``no sé''.\\
        \textit{[\ldots 2.000 tokens más\ldots]}
      \end{quote}

      \vspace{0.1em}
      {\color{fenred}\textbf{Riesgo:}} la instrucción crítica queda diluida. El modelo puede ignorarla.
    \end{column}
    \begin{column}{0.48\textwidth}
      {\color{fenblue}\textbf{Prompt bien estructurado:}}
      \vspace{0.2em}

      \begin{quote}\small
        \textbf{Sistema:} ``Si no tienes certeza, responde: 'No tengo información suficiente'. Usa solo la normativa adjunta.''\\
        \textit{[normativa bien segmentada]}
      \end{quote}

      \vspace{0.1em}
      {\color{fenblue}\textbf{Resultado:}} la instrucción recibe alta atención; el modelo la respeta.
    \end{column}
  \end{columns}

\end{frame}

% ----------------------------------------------------------

\begin{frame}{Lo que la atención NO garantiza}

  \begin{columns}[T]
    \begin{column}{0.5\textwidth}
      A pesar de su potencia, la atención \textbf{no resuelve todo}:
      \vspace{0.4em}
      \begin{itemize}
        \item \highlight{No garantiza verdad factual} --- puede ponderar bien tokens incorrectos
        \item \highlight{No garantiza razonamiento correcto} --- pesar bien $\neq$ razonar bien
        \item \highlight{No garantiza priorización adecuada} --- depende de cómo está construido el contexto
        \item \highlight{No elimina las alucinaciones}
      \end{itemize}
    \end{column}
    \begin{column}{0.47\textwidth}
      \begin{alertblock}{Regla simple}
        Si el contexto está mal armado, el modelo puede enfocarse en lo menos importante.\\[0.4em]
        \textbf{La calidad del input determina la calidad de la atención.}
      \end{alertblock}
      \vspace{0.3em}
      Por eso la gobernanza y el diseño cuidadoso del prompt son parte de cualquier solución seria.
    \end{column}
  \end{columns}

\end{frame}

% ===========================================================
\section{Síntesis}
% ===========================================================

\begin{frame}{Síntesis: ¿qué cambió Attention Is All You Need?}

  \begin{center}
  \begin{tabular}{>{\bfseries}p{3.5cm} p{8cm}}
    \toprule
    Dimensión & Aporte \\
    \midrule
    Arquitectura & Reemplazó RNN/LSTM por atención pura \\
    Paralelización & Permitió entrenar en paralelo $\rightarrow$ escalar a miles de millones de parámetros \\
    Contexto global & Cualquier token puede relacionarse con cualquier otro directamente \\
    Fundamento & Es la base de GPT, BERT, Claude, Gemini, Llama y prácticamente todos los LLMs vigentes \\
    Práctica & La calidad del prompt influye en qué partes del contexto reciben más atención \\
    \bottomrule
  \end{tabular}
  \end{center}

  \vspace{0.5em}
  \begin{center}
    \large \textit{``Gran parte del salto de los LLMs modernos pasa por este mecanismo.''}
  \end{center}

\end{frame}

% ----------------------------------------------------------

\begin{frame}{Actividad de reflexión}

  \begin{block}{Para discutir en grupos (15--20 min)}
    Con base en el abstract e introducción del paper (Vaswani et al., 2017):
    \begin{enumerate}
      \item ¿Qué limitación de los modelos anteriores identifica el paper?
      \item ¿Cuál es la propuesta central del Transformer?
      \item ¿Por qué la atención aparece como una innovación tan importante?
      \item ¿Qué conexión ven entre este paper y los LLMs que usamos hoy?
    \end{enumerate}
  \end{block}

  \vspace{0.4em}
  \textbf{Producto esperado --- mini-síntesis de 3 partes:}
  \begin{description}
    \item[Problema anterior] ¿qué estaba fallando o limitando a los modelos previos?
    \item[Aporte del paper] ¿qué cambia con el Transformer?
    \item[Impacto actual] ¿por qué esto sigue importando hoy?
  \end{description}

\end{frame}

% ----------------------------------------------------------

\begin{frame}[plain]
  \begin{center}
    \vspace{2em}
    {\Large \textbf{Gracias}}\\[1em]
    {\normalsize Sebastián Egaña Santibáñez}\\[0.3em]
    {\small Módulo 3 · IA Generativa y LLMs}\\
    {\small FEN Unegocios, Universidad de Chile}\\[1.5em]
    \textcolor{fengray}{\small Vaswani, A. et al. (2017). \textit{Attention Is All You Need}. NeurIPS.}\\
    \textcolor{fengray}{\small \texttt{https://arxiv.org/abs/1706.03762}}
  \end{center}
\end{frame}

% ============================================================
\end{document}
% ============================================================
